\bookchapter{บทที่ 51: พัฒนาหน้าจอ — เริ่มต้นกับ Jetpack Compose}

\chapterintro{หน้าจอคือทุกสิ่งที่ผู้ใช้เห็น ก่อน Compose จะเกิดขึ้น หน้าจอ Android เขียนด้วย XML — ยุ่งยากเหมือนทำโจทย์เติมช่องว่าง; ปี 2021 Google เปิดตัว Jetpack Compose: เขียนโค้ด Kotlin “วาด” หน้าจอไปตรง ๆ เห็นอย่างไรได้อย่างนั้น แก้โค้ดหนึ่งบรรทัดพรีวิวก็เด้งใหม่ทันที มาถึงปี 2026 Compose กลายเป็นกระแสหลักของโปรเจกต์ใหม่เด็ดขาด ในบทนี้ เราจะเรียน “วาดรูป” กัน}

\booksection{แนวคิดของ Compose: หน้าจอก็คือฟังก์ชัน}

แนวคิดหลักของ Compose พูดในประโยคเดียว: หน้าจอ = ฟังก์ชัน ทุกอย่างที่แสดงบนหน้าจอได้ (ข้อความ ปุ่ม รูปภาพ) ล้วนเป็น “ฟังก์ชันประกอบได้” (Composable) ที่มาร์กด้วย @Composable เอามาซ้อนกันเหมือนต่อบล็อกเลโก้ ก็กลายเป็นหน้าจอ

\begin{lstlisting}[language=Kotlin]
@Composable
fun Greeting(name: String) {
    Text(text = "สวัสดี $name!")     // บล็อก "ข้อความ" หนึ่งชิ้น
}

// ประกอบซ้อนกัน: ชิ้นใหญ่ครอบชิ้นเล็ก
@Composable
fun HomeScreen() {
    Column {                          // เรียงตามแนวตั้ง
        Greeting(name = "ซานผวน")
        Greeting(name = "ลี่ซือ")
        Button(onClick = { /* แตะแล้วทำอะไร */ }) {
            Text("แตะฉัน")
        }
    }
}
\end{lstlisting}

สังเกต “หมวก” ที่ชื่อ @Composable — มันบอกคอมไพเลอร์ว่า “นี่คือบล็อกต่อหน้าจอ” บล็อกซ้อนกันได้ไม่จำกัด นี่คือความลับทั้งหมดของ Compose

\booksection{คอมโพเนนต์พื้นฐาน: ข้อความ ปุ่ม ช่องกรอก รูปภาพ}

คอมโพเนนต์ของ Compose ก็คือ “บล็อกต่อ” ที่เจอบ่อยที่สุดมีสี่ตัว:

\begin{center}
\begin{tabularx}{\linewidth}{|X|X|X|}
\hline
\textbf{คอมโพเนนต์} & \textbf{ทำอะไร} & \textbf{พารามิเตอร์ที่ต้องมี} \tabularnewline
\hline
Text & แสดงข้อความ & text (ข้อความที่จะแสดง) \tabularnewline
Button & ปุ่ม & onClick (ฟังก์ชันเรียกกลับตอนกด) \tabularnewline
TextField & ช่องกรอกข้อความ & value, onValueChange (ฟังก์ชันเรียกกลับเมื่อข้อความเปลี่ยน) \tabularnewline
Image & แสดงรูปภาพ & painter (ทรัพยากรรูปภาพ) \tabularnewline
\hline
\end{tabularx}
\end{center}

\begin{center}\small\color{gray}คอมโพเนนต์พื้นฐานหลักสี่ตัวของ Compose\end{center}

\begin{lstlisting}[language=Kotlin]
@Composable
fun BasicWidgets() {
    Column {
        Text(text = "นี่คือข้อความ")
        Button(onClick = { println("ถูกแตะแล้ว") }) {
            Text("ปุ่ม")
        }
        var input by remember { mutableStateOf("") }
        TextField(
            value = input,
            onValueChange = { input = it },
            label = { Text("พิมพ์อะไรก็ได้") }
        )
    }
}
\end{lstlisting}

ในสองบรรทัดของช่องกรอกนั้นปรากฏ remember \{ mutableStateOf(“”) \} — นี่คือเวทมนตร์เรื่อง “สถานะ” ของ Compose จะอธิบายละเอียดในหัวข้อที่ 5 ของบทนี้ ตอนนี้แค่ทำหน้าคลุกเคล้าก่อน

\booksection{สามชิ้นเลย์เอาต์: Column, Row, Box}

มีคอมโพเนนต์แล้ว จัดวางยังไง? จำ “กล่อง” สามใบนี้:

\begin{itemize}
  \item Column: เรียงแนวตั้ง (เหมือนต่อแถวรับอาหาร);
  \item Row: เรียงแนวนอน (เหมือนนั่งเคียงกัน);
  \item Box: วางซ้อนกัน (เหมือนไพ่กองทับกัน เหมาะทำป้ายมุม เลเยอร์ซ้อน)
\end{itemize}

\begin{lstlisting}[language=Kotlin]
Column {
    Text("บรรทัดแรก")
    Row {
        Text("ซ้าย")
        Text("ขวา")
    }
    Box {
        Text("ชั้นล่าง")
        Text("ป้ายเล็กที่ซ้อนอยู่ด้านบน")
    }
}
\end{lstlisting}

อยากให้เลย์เอาต์สวยขึ้น? เพิ่ม Modifier (ตัวแต่งแต้ม) ให้คอมโพเนนต์ — เหมือนแต่งหน้าให้บล็อก:

\begin{lstlisting}[language=Kotlin]
Text(
    text = "ข้อความสวย ๆ",
    modifier = Modifier
        .padding(16.dp)          // เว้นรอบด้าน
        .fillMaxWidth()          // กว้างเต็มแถว
        .background(Color.Blue)  // สีพื้นหลัง
        .clickable { /* แตะ */ }// แตะได้
)
// ลูกโซ่: เรียงต่อกันทีละตัว ลำดับมีผล (padding ก่อน แล้วพื้นหลัง พื้นหลังไม่รวมช่องว่าง)
\end{lstlisting}

\begin{tcbwarn}{ข้อควรรู้: ลำดับของ Modifier มีผล}
padding(16.dp).background(Blue) กับ background(Blue).padding(16.dp) ผลลัพธ์ต่างกัน — แบบแรกคือ “เว้นระยะก่อนแล้วค่อยราบสี” (พื้นหลังไม่รวมขอบ) ส่วนแบบหลังคือ “ราบสีก่อนแล้วค่อยเว้นระยะ” (พื้นหลังรวมขอบ) มือใหม่เจอ “ขอบขาวประหลาด” ที่วาดผิดตรงนี้บ่อยที่สุด ถ้าไม่แน่ใจ ลองสลับลำดับดู พรีวิวจะบอกคำตอบทันที

\end{tcbwarn}

\booksection{สถานะ: ทำให้หน้าจอ “มีชีวิต”}

หน้าจอนิ่ง ๆ ไม่สนุก อยากให้กดปุ่มแล้วเลข +1 ให้พิมพ์ในช่องกรอกได้ ก็ต้องมีสถานะ (State) กฎของ Compose: หน้าจอสะท้อนสถานะอย่างซื่อสัตย์เสมอ สถานะเปลี่ยนเมื่อไร หน้าจอวาดใหม่ให้เอง

\begin{lstlisting}[language=Kotlin]
@Composable
fun Counter() {
    // remember + mutableStateOf: จำค่าที่เปลี่ยนได้นี้ไว้
    var count by remember { mutableStateOf(0) }

    Column {
        Text(text = "คุณแตะไป $count ครั้ง")
        Button(onClick = { count++ }) {   // แก้สเตต
            Text("แตะฉัน +1")                // หน้าจอเด้งเอง
        }
    }
}
\end{lstlisting}

สามบรรทัด ได้ตัวนับที่มีชีวิต นี่คือจุดที่ Compose สบายกว่าสไตล์เก่า: ไม่ต้อง “รีเฟรชหน้าจอ” ด้วยมือเอง แค่แก้ข้อมูลก็พอ จำสองฟังก์ชันสำคัญไว้: remember (ทำให้ค่าไม่ถูกล้างตอนหน้าจอ recompose) กับ mutableStateOf (บอก Compose ว่า “ค่านี้เปลี่ยนแล้วต้องวาดใหม่”)

ผู้ใช้แตะ

สเตตเปลี่ยน (count++)

Compose ตรวจจับการเปลี่ยนแปลง

วาด UI ใหม่อัตโนมัติ (recomposition)

\begin{center}\small\color{gray}ภาพที่ 3: การไหลของข้อมูลทางเดียว “สถานะ หน้าจอ” ของ Compose\end{center}

\booksection{ธีมกับโหมดมืด: ตั้งครั้งเดียว มีผลทั้งแอป}

MyFirstAppTheme \{ ... \} ในโปรเจกต์เทมเพลตนั่นแหละคือธีม มันอิงมาตรฐานดีไซน์ Material 3 ช่วยจัดการสี ฟอนต์ รูปทรงให้เป็นระบบเดียวกัน รองรับโหมดมืดแค่สามขั้นตอน:

\begin{enumerate}
  \item นิยามชุดสีแบบสว่าง/มืดไว้ในธีม (เทมเพลตสร้าง ColorScheme ให้เองอยู่แล้ว);
  \item เมื่อระบบสลับไปโหมดมืด isSystemInDarkTheme() จะตัดสินใจให้เอง;
  \item ใช้สีจากธีม (เช่น MaterialTheme.colorScheme.primary) แทนการเขียนค่าสีฝังตาย
\end{enumerate}

\begin{tcbinfo}{เคล็ดลับ: สีแบบไดนามิกของ Material 3}
ตั้งแต่ Android 12 เป็นต้นมา “Material You” ทำให้แอปดึงสีจากวอลเปเปอร์ของคุณได้ เทมเพลต Compose รองรับตามค่าเริ่มต้น (dynamicColor) แอปของคุณจะ “เปลี่ยนชุด” เองบนมือถือผู้ใช้แต่ละคน นี่คือความโรแมนติกแบบ Android ยุคใหม่

\end{tcbinfo}

\booksection{พรีวิว อนิเมชัน และการแก้สด: “ไอเทมเสริมพลัง” ของ Compose}

ของเสริมพลังสามอย่างเรื่องประสบการณ์พัฒนาของ Compose พอใช้แล้วย้อนกลับไปแบบเดิมไม่ได้อีกเลย:

\begin{itemize}
  \item พรีวิว @Preview: ใส่ annotation @Preview ให้ Composable Android Studio จะเรนเดอร์หน้าจอทางขวาทันที — แก้โค้ดเห็นผลทันทีทันใด ไม่ต้องติดตั้งลงมือถือทุกครั้ง;
  \item การแก้สด (Live Edit): Android Studio ปี 2026 รองรับ "อัปเดตร้อน" การแก้บางส่วนบนโปรแกรมจำลอง/เครื่องจริง ปรับเล็ก ๆ อย่างเปลี่ยนสีข้อความ เห็นผลบนเครื่องภายในวินาทีเดียว;
  \item API อนิเมชัน: animateFloatAsState, AnimatedVisibility ใช้แค่ไม่กี่บรรทัดก็ให้เอฟเฟกต์เนียน ๆ — "ความพรีเมียม" มักขาดแค่สิบกว่าบรรทัดแบบนี้เอง
\end{itemize}

\begin{lstlisting}[language=Kotlin]
@Composable
fun AnimatedBox() {
    var big by remember { mutableStateOf(false) }
    // เปลี่ยนค่าแล้วเลื่อนไปค่าใหม่อย่างนุ่มนวลอัตโนมัติ
    val size by animateDpAsState(if (big) 120.dp else 60.dp)

    Box(Modifier
        .size(size)
        .background(MaterialTheme.colorScheme.primary)
        .clickable { big = !big }
    )
}
\end{lstlisting}

\begin{tcbinfo}{เคล็ดลับ: ฟังก์ชันก่อน อนิเมชันทีหลัง}
โรคประจำตัวของมือใหม่คือ “หน้าจอยังไม่ถูก ก็ไปจุกจิกกับอนิเมชัน” ให้ฟังก์ชันรันผ่านก่อน แล้วค่อยใส่เอฟเฟกต์ทีเดียวท้ายสุด — นี่ก็คือหลัก “เรียบง่ายเหนือสิ่งใด”: อนิเมชันคือดอกไม้ประดับ ฟังก์ชันคือถ่านยามหนาว

\end{tcbinfo}

\begin{tcbpractice}{ลงมือทำด้วยกัน: สร้าง "ปุ่มอารมณ์"}
แปลงโปรเจกต์เทมเพลตเป็นเครื่องสลับอารมณ์: ปุ่มเดียว กดแล้ววนสลับระหว่าง ดีใจ เศร้า โมโห ข้อความกับสีพื้นหลังเปลี่ยนตาม คำใบ้: ใช้ when เลือกข้อความกับสีตามสถานะ เสร็จแล้วส่งให้เพื่อนเล่น เก็บกระแสคลื่นแรกว่า “โห ทำแอปเป็นแล้ว” ดู

\end{tcbpractice}

\begin{tcbsummary}{สรุปบทนี้}
\begin{itemize}
  \item แนวคิด Compose: หน้าจอ = ฟังก์ชัน @Composable ซ้อนกันแบบต่อบล็อก
  \item คอมโพเนนต์หลักสี่ตัว: Text / Button / TextField / Image
  \item เลย์เอาต์สามชิ้น: Column เรียงแนวตั้ง, Row เรียงแนวนอน, Box วางซ้อน
  \item Modifier แต่งตัวแบบลูกโซ่: padding, fillMaxWidth, background, clickable — ลำดับมีความหมาย
  \item สถานะคือดวงวิญญาณ: remember + mutableStateOf แก้ข้อมูลเมื่อไรหน้าจอวาดใหม่เอง
  \item ธีมใช้ Material 3 + สีไดนามิก โหมดมืดตั้งครั้งเดียวมีผลทั้งแอป
\end{itemize}

\end{tcbsummary}
